%% LyX 2.4.4 created this file.  For more info, see https://www.lyx.org/.
%% Do not edit unless you really know what you are doing.
\documentclass[english,t]{beamer}
\usepackage{mathptmx}
\usepackage[T1]{fontenc}
\usepackage[latin9]{inputenc}
\usepackage{amssymb}
\usepackage{cancel}

\makeatletter
%%%%%%%%%%%%%%%%%%%%%%%%%%%%%% Textclass specific LaTeX commands.
% this default might be overridden by plain title style
\newcommand\makebeamertitle{\frame{\maketitle}}%
% (ERT) argument for the TOC
\AtBeginDocument{%
  \let\origtableofcontents=\tableofcontents
  \def\tableofcontents{\@ifnextchar[{\origtableofcontents}{\gobbletableofcontents}}
  \def\gobbletableofcontents#1{\origtableofcontents}
}

%%%%%%%%%%%%%%%%%%%%%%%%%%%%%% User specified LaTeX commands.
\usetheme{Warsaw}
% or ...

\setbeamercovered{transparent}
% or whatever (possibly just delete it)
\setbeamercovered{invisible} %makes overlays invisible

%gets rid of navigation symbols
\setbeamertemplate{navigation symbols}{}

%\usepackage{hyperref}
%\usepackage[usenames,dvipsnames,table]{xcolor} %adds additional color options
\definecolor{links}{HTML}{2A1B81}
%\hypersetup{colorlinks,linkcolor=blue,urlcolor=links}

\usepackage{multicol}

\usepackage{tikz}
\usetikzlibrary{calc}
\usetikzlibrary{plotmarks}
\usetikzlibrary{shapes}
\usepackage{pgfplots}
\pgfplotsset{compat=newest}

\usepackage{forest} %%for tree diagram

\usepackage{calc}
\usepackage[nomessages]{fp}

%\usepackage{advdate}    % Advancing/saving dates
\usepackage{datetime}   % Dates formatting
%\usepackage{datenumber} % Counters for dates

\usepackage[super]{nth}

\usepackage{etex}

\usepackage[outline]{contour}  %allows text halos
\contourlength{1pt} %sets halo width
%%example: \node at (0.75,0.5) {\contour{white}{\Large Text!}};
%%example:  \node [fill=white,inner sep=1pt] at (2.25,0.5) {\Large 			Text!};
%%example:  \node [fill=white,rounded corners=2pt,inner sep=1pt] at 			(3.75,0.5) {\Large Text!};

%%%%%%%%%%%%%%%%%%%%%%%%%%%
\def\formatdateny{\csname noyear\languagename\expandafter\endcsname\formatdate}
\def\noyearenglish#1, \the\@year{#1}
\let\noyearamerican\noyearenglish
\let\noyearbritish\noyearenglish

%%allows uncover with forest diagrams
\tikzset{
    invisible/.style={opacity=0,text opacity=0},
    visible on/.style={alt=#1{}{invisible}},
    alt/.code args={<#1>#2#3}{%
      \alt<#1>{\pgfkeysalso{#2}}{\pgfkeysalso{#3}} % \pgfkeysalso doesn't change the path
    },
}
\forestset{
  visible on/.style={
    for tree={
      /tikz/visible on={#1},
      edge={/tikz/visible on={#1}}}}}

%%%redefines column alignment to match vertically
%\define@key{beamerframe}{t}[true]{% top
  %\beamer@frametopskip=.2cm plus .5\paperheight\relax%
  %\beamer@framebottomskip=0pt plus 1fill\relax%
 % \beamer@frametopskipautobreak=\beamer@frametopskip\relax%
  %\beamer@framebottomskipautobreak=\beamer@framebottomskip\relax%
  %\def\beamer@initfirstlineunskip{}%
%}

%%provides pdf with 4 slides per page; don't forget to add 'handout' to remove overlays
%\usepackage{pgfpages}
%\pgfpagesuselayout{4 on 1}[a4paper, landscape, border shrink=7mm]
%\pgfpageslogicalpageoptions{1}{border code=\pgfusepath{stroke}}
%\pgfpageslogicalpageoptions{2}{border code=\pgfusepath{stroke}}
%\pgfpageslogicalpageoptions{3}{border code=\pgfusepath{stroke}}
%\pgfpageslogicalpageoptions{4}{border code=\pgfusepath{stroke}}
%%%Don't forget to add 'handout'

\makeatother

\usepackage{babel}
\begin{document}
\newcommand{\xmax}{2000}
\newcommand{\xmin}{0}
\newcommand{\xstep}{0} %must be xmin+n
\newcommand{\ymax}{500}
\newcommand{\ymin}{0}
\newcommand{\ystep}{0} %must be ymin+n

\newcounter{countEx} \setcounter{countEx}{0}
\newcounter{countProb} \setcounter{countProb}{0}
\resetcounteronoverlays{countEx} \resetcounteronoverlays{countProb}

\newcommand{\ex}[3]{
	\addtocounter{countEx}{1}
	\only<#1-#2>{\begin{exampleblock} {Example \chNum.\secNum.\arabic{countEx}.} #3	\end{exampleblock}}
}

\newcommand{\prob}[4]{
	\addtocounter{countProb}{1}
	\only<#1-#2>{\begin{exampleblock} {Problem  \chNum.\secNum.\arabic{countProb}.\,#3} #4	\end{exampleblock}}
}

\newcommand{\yline}[4]{
	(#2 - #4)/(#1 - #3)*(x - #1) + #2
}

\beamerdefaultoverlayspecification{<+->}

%%%%Chapter and Section numbers%%%%%
\newcommand{\chNum}{2}
\newcommand{\secNum}{3}
\newcommand{\secTitle}{Introducion to Combinatorics}

\title[Section \chNum.\secNum: \secTitle]{Section \chNum.\secNum: \secTitle}

\subtitle{Finite Mathematics~~MTH 107~~Fall 2014}

\author{Dr.\,James Ashe}

\titlegraphic{\includegraphics[height=2.75cm]{/home/jim/JamesAshe/MHU/images/MarsHilUniversityLogo-Virtical.png}}

\date{\formatdate{27}{8}{2014}}

\institute{Mars Hill University}

\makebeamertitle 
\begin{frame}{}

\begin{exampleblock}<1->{}
How many different outfits could you create from two pants, three
shirts, and a pair of shoes?
\end{exampleblock}
\begin{itemize}
\item <4->\emph{Fundamentals of Counting.}
\end{itemize}
\begin{exampleblock}<2->{}
A pizza place lets you choose any $3$ of $10$ different toppings
for your pizza. How many different pizzas could you create?
\end{exampleblock}
\begin{itemize}
\item <5->\emph{Combinations. }
\item []<6->Order doesn't matter.
\end{itemize}
\begin{exampleblock}<3->{}
A baseball team has $9$ players. How many different line-ups are
possible?
\end{exampleblock}
\begin{itemize}
\item <7->\emph{Permutations.}
\item []<8->Order does matter.
\end{itemize}
\end{frame}

\begin{frame}{Fundamental Principle of Counting}

\ex{1}{}{How many different outfits can worn from $3$ pants, $2$ shirts, and $1$ pair of shoes?}

\forestset{  
	L1/.style={rectangle, rounded corners, shade, top color=white,     bottom color=blue!50!black!20, draw=blue!40!black!60, very     thick},   	
	L2/.style={rectangle, rounded corners, shade, top color=white,     bottom color=blue!50!black!20, draw=blue!40!black!60, very     thick,
		edge={black,line width=1.15pt,line cap=round}},   	
	L3/.style={rectangle, rounded corners, shade, top color=white,     bottom color=blue!50!black!20, draw=blue!40!black!60, very     thick,
		edge={black,line width=1.15pt,line cap=round}},   	
	L4/.style={rectangle, rounded corners, shade, top color=white,     bottom color=blue!50!black!20, draw=blue!40!black!60, very     thick,
	edge={black,line width=1.15pt,line cap=round}}, 
}

\begin{forest}     
	for tree={
		scale=.75,        
		grow=0,
		reversed, %tree direction         
		parent anchor=east,
		child anchor=west, %edge anchors   		edge={line cap=round},
		outer sep=+1pt, %edge/node connection  		rounded corners,
		minimum width=15mm,
		minimum height=8mm, %node shape         l 		sep=10mm %level distance     
		}   
[Start,L1,for children={visible on=<2->}
	[Pants 1,L2, for children={visible on=<3->}         
		[Shirt 1,L3,for children={visible on=<6->}             
			[Shoes 1,L4,]        
		]         
		[Shirt 2,L3,for children={visible on=<6->}            
			[Shoes 1,L4]
		]               
	]     
	[Pants 2,L2,for children={visible on=<4->}         
		[Shirt 1,L3,for children={visible on=<6->}            
			[Shoes 1,L4]         
		]         
		[Shirt 2,L3,for children={visible on=<6->}           
			[Shoes 1,L4]
		]               
	] 
	[Pants 3,L2,for children={visible on=<5->}         
		[Shirt 1,L3,for children={visible on=<6->}            
			[Shoes 1,L4]         
		]         
		[Shirt 2,L3,for children={visible on=<6->}           
			[Shoes 1,L4]
		]               
	] 
] 
\end{forest}
\end{frame}

\begin{frame}{Fundamental Principle of Counting}

\begin{block}{Fundamental Priniple of Counting}
If there are $a$ possible outcomes for an event and $b$ number
of possible outcomes for another event, then there are $a\times b$
possible outcomes for both events together.
\end{block}
\begin{itemize}
\item <2->[1.]Write blanks for each event/decision.
\item <4->[2.]Enter the number of outcomes/choices for each event/decision
and multiply. 
\end{itemize}
\ex{1}{}{How many different outfits can worn from $3$ pants, $4$ shirts, and $10$ pairs of shoes?}
\begin{itemize}
\item [] \only<3-4>{$\underline{\,\,\,\,}\times\underline{\,\,\,\,}\times\underline{\,\,\,\,}$}\only<5-5>{$\underline{3}\times\underline{\,\,\,\,}\times\underline{\,\,\,\,}$}\only<6-6>{$\underline{3}\times\underline{4}\times\underline{\,\,\,\,}$}\only<7-7>{$\underline{3}\times\underline{4}\times\underline{10}$}\only<8-8>{$\underline{3}\times\underline{4}\times\underline{10}=120$}
\end{itemize}
\end{frame}

\begin{frame}{Fundamental Principle of Counting}

\ex{1}{}{Three coins are tossed.}
\begin{itemize}
\item <2->[a.]Use the Fundamental Principle of Counting to determine how
many different outcomes are possible. 
\item <3->[b.]Construct a tree diagram to list all possible outcomes.
\end{itemize}
\end{frame}

\begin{frame}{Factorials}

\begin{itemize}
\item []<1->How different ways can we arrange $3$ things?
\item []
\item []<2->How different ways can we arrange $4$ things?
\item []
\item []<3->How different ways can we arrange $5$ things?
\item []
\item []<4->How different ways can we arrange $n$ things?
\end{itemize}
\end{frame}

\begin{frame}{Factorials}

\begin{columns}


\column{5.5cm}
\begin{block}{Factorial}
Let $n$ be a positive integer, then $n$ \emph{factorial }is

\begin{itemize}
\item []$n!=n\cdot(n-1)\cdot(n-2)\dots\cdot2\cdot1$
\item []<2->Special case: $0!=1$.
\end{itemize}
\end{block}

\column{5.5cm}
\begin{itemize}
\item <3->[ex.]$5!$\only<4->{$=5\cdot4\cdot3\cdot2\cdot1=120$}
\item []
\item <5->[ex.]$\frac{7!}{5!}$\only<6->{
\[
=\frac{7\cdot6\cdot5\cdot4\cdot3\cdot2\cdot1}{\,\,\,\,\,\,\,\,\,\,\,\,5\cdot4\cdot3\cdot2\cdot1}
\]
}\only<7->{$=42$}
\item <8->[ex.]
\[
\frac{7!}{3!\cdot5!}
\]
\only<9->{
\[
=\frac{\,\,\,\,\,7\cdot6\cdot5\cdot4\cdot3\cdot2\cdot1}{(3\cdot2\cdot1)(5\cdot4\cdot3\cdot2\cdot1)}
\]
}\only<10->{$=\frac{7\cdot6}{6}=7$}
\end{itemize}
\end{columns}
\end{frame}

\begin{frame}{Problems}

\prob{1}{1}{A sandwhich shop has $2$ different breads, $8$ meats, $5$ cheeses, and $4$ sauces.}{How many different sandwhiches are possible?}

\prob{2}{3}{A password has $2$ letters   followed by $4$ digits (0-9).}{How many different passwords are possible if,
\begin{itemize} 
	\item <2->[a.]Letters and digits can be repeated. 
	\item <3->[b.]Letters and digits cannot be repeated.
\end{itemize}
}

\prob{4}{}{Find the indicated value}{
\[
\frac{90!}{89!}\hspace{1em}\hspace{1em}\hspace{1em}\hspace{1em}\hspace{1em}\hspace{1em}\hspace{1em}\hspace{1em}\hspace{1em}\hspace{1em}\hspace{1em}\hspace{1em}\hspace{1em}\hspace{1em}\hspace{1em}\hspace{1em}
\]
}\vspace{1.25cm}\prob{5}{}{Find the indicated value}{
\[
\frac{90!}{88!\cdot3!}\hspace{1em}\hspace{1em}\hspace{1em}\hspace{1em}\hspace{1em}\hspace{1em}\hspace{1em}\hspace{1em}\hspace{1em}\hspace{1em}\hspace{1em}\hspace{1em}\hspace{1em}\hspace{1em}\hspace{1em}\hspace{1em}
\]
}
\end{frame}

\end{document}
